\documentclass{article}
\usepackage{amsmath,amssymb}
\usepackage[margin=1in]{geometry}

\begin{document}


\begin{flushright}
Neil Hwang\\
Linear Algebra\\
March 25, 2026\\
Homework 1\\

\end{flushright}





\bigskip
\begin{center}
\rule{0.5\textwidth}{1pt}\\[6pt]
\textbf{Homework from Lecture on February 3, 2026}\\[2pt]
\rule{0.5\textwidth}{1pt}
\end{center}
\bigskip

\clearpage
\hfill \break  $~\!\!$ \dotfill \medskip \\
\textbf{Question.}\\
Let $W \subset V$. Show that $W$ is a subspace of $V$ if and only if:
\begin{align*}
\forall x,y \in W,\quad x+y \in W, \\
\forall a \in F,\ \forall x \in W,\quad ax \in W, \\
0 \in W.
\end{align*}

\noindent\\
\textbf{Solution.}\\
\quad\\
\noindent\fbox{ \parbox{\textwidth}{
Suppose first that $W$ is a subspace of $V$. Then $W$ is itself a vector space with the same operations as $V$. Hence:
\begin{align*}
x,y\in W &\implies x+y\in W, \\
a\in F,\ x\in W &\implies ax\in W, \\
0&\in W.
\end{align*}

Conversely, assume that $W\subset V$ satisfies
\begin{align*}
x,y\in W &\implies x+y\in W, \\
a\in F,\ x\in W &\implies ax\in W, \\
0&\in W.
\end{align*}
We show that $W$ is a subspace of $V$, i.e., $W$ is itself a vector space under the same operations as $V$. We verify all eight vector space axioms.

Since $W \subset V$, addition and scalar multiplication on $W$ are the restrictions of the operations on $V$. All elements of $W$ are also elements of $V$, so any identity that holds for all elements of $V$ automatically holds for elements of $W$.

\textbf{(A1) Commutativity of addition:} For all $x, y \in W$, since $x, y \in V$ and $V$ is a vector space,
\[
x + y = y + x.
\]

\textbf{(A2) Associativity of addition:} For all $x, y, z \in W \subset V$,
\[
(x + y) + z = x + (y + z).
\]

\textbf{(A3) Existence of additive identity:} $0 \in W$ by assumption, and $x + 0 = x$ for all $x \in W$ since this holds in $V$.

\textbf{(A4) Existence of additive inverses:} Let $x \in W$. Since $-1 \in F$ and $W$ is closed under scalar multiplication,
\[
-x = (-1)x \in W.
\]
And $x + (-x) = 0$ holds because it holds in $V$.

\textbf{(S1) Compatibility of scalar multiplication:} For all $a, b \in F$ and $x \in W \subset V$,
\[
a(bx) = (ab)x.
\]

\textbf{(S2) Identity element of scalar multiplication:} For all $x \in W \subset V$,
\[
1 \cdot x = x.
\]

\textbf{(D1) Distributivity of scalar multiplication over vector addition:} For all $a \in F$ and $x, y \in W \subset V$,
\[
a(x + y) = ax + ay.
\]

\textbf{(D2) Distributivity of scalar multiplication over field addition:} For all $a, b \in F$ and $x \in W \subset V$,
\[
(a + b)x = ax + bx.
\]

In each case, the identity holds because $W \subset V$ and the identity already holds in $V$. The only axioms that require the three closure assumptions are: (A3) requires $0 \in W$; (A4) requires closure under scalar multiplication (to get $-x = (-1)x \in W$); and all axioms involving addition or scalar multiplication of elements of $W$ require that the results remain in $W$ (closure under addition and scalar multiplication).

Therefore all vector space axioms hold on $W$, so $W$ is a subspace of $V$.
}}

\bigskip


\clearpage
\hfill \break  $~\!\!$ \dotfill \medskip \\
\textbf{Question.}\\
Let $T: V \to W$ be a function. Show that $T$ is linear if and only if:
\begin{align*}
T(x+y) = T(x) + T(y), \\
T(ax) = aT(x)
\end{align*}
for all $x,y \in V$ and $a \in F$.

\noindent\\
\textbf{Solution.}\\
\quad\\
\noindent\fbox{ \parbox{\textwidth}{
By definition, $T$ is linear if for all $x,y\in V$ and all $a,b\in F$,
\begin{align*}
T(ax+by)=aT(x)+bT(y).
\end{align*}

Suppose first that $T$ is linear. Then taking $a=b=1$, we get
\begin{align*}
T(x+y)=T(1x+1y)=1T(x)+1T(y)=T(x)+T(y).
\end{align*}
Taking $b=0$ and $y=0$, we get
\begin{align*}
T(ax)=T(ax+0\cdot 0)=aT(x)+0\cdot T(0)=aT(x).
\end{align*}
So the two displayed properties follow.

Conversely, suppose that for all $x,y\in V$ and all $a\in F$,
\begin{align*}
T(x+y)=T(x)+T(y),\qquad T(ax)=aT(x).
\end{align*}
We show that $T$ is linear. Let $x,y\in V$ and $a,b\in F$. Then
\begin{align*}
T(ax+by)
&=T(ax)+T(by) \\
&=aT(x)+bT(y).
\end{align*}
Thus $T$ satisfies the defining property of linearity, hence $T$ is linear.
}}

\bigskip

\clearpage
\hfill \break  $~\!\!$ \dotfill \medskip \\
\textbf{Question.}\\
Let $B = \{x_1,\dots,x_n\}$ be a basis of $V$. Show that every $x \in V$ can be written uniquely as:
\begin{align*}
x = a_1 x_1 + \cdots + a_n x_n.
\end{align*}

\noindent\\
\textbf{Solution.}\\
\quad\\
\noindent\fbox{ \parbox{\textwidth}{
Since $B$ is a basis of $V$, by definition it spans $V$ and is linearly independent. Because $B$ spans $V$, every $x\in V$ can be expressed as a linear combination of the basis vectors:
\begin{align*}
x=a_1x_1+\cdots+a_nx_n
\end{align*}
for some scalars $a_1,\dots,a_n\in F$. This proves existence.

To prove uniqueness, suppose that
\begin{align*}
x=a_1x_1+\cdots+a_nx_n=b_1x_1+\cdots+b_nx_n.
\end{align*}
Subtracting the two expressions gives
\begin{align*}
0=(a_1-b_1)x_1+\cdots+(a_n-b_n)x_n.
\end{align*}
Since $B$ is linearly independent, the only linear combination of $x_1,\dots,x_n$ equal to $0$ is the trivial one. Therefore
\begin{align*}
a_1-b_1=0,\quad \dots,\quad a_n-b_n=0.
\end{align*}
Hence
\begin{align*}
a_i=b_i\qquad \text{for all } i=1,\dots,n.
\end{align*}
So the representation is unique.
}}

\bigskip

\clearpage
\hfill \break  $~\!\!$ \dotfill \medskip \\
\textbf{Question.}\\
Let $V$ be an $n$-dimensional vector space with ordered basis $B = \{x_1,\dots,x_n\}$. Define the coordinate map:
\begin{align*}
[x]_B = (a_1,\dots,a_n) \in F^n
\end{align*}
where $x = a_1 x_1 + \cdots + a_n x_n$.

Show that the map $x \mapsto [x]_B$ is a linear isomorphism from $V$ to $F^n$.

\noindent\\
\textbf{Solution.}\\
\quad\\
\noindent\fbox{ \parbox{\textwidth}{
Define $\Phi:V\to F^n$ by
\begin{align*}
\Phi(x)=[x]_B,
\end{align*}
where if $x=a_1x_1+\cdots+a_nx_n,$ then $[x]_B=(a_1,\dots,a_n).$

We first show that $\Phi$ is linear. Let
\begin{align*}
x=a_1x_1+\cdots+a_nx_n,\qquad y=b_1x_1+\cdots+b_nx_n.
\end{align*}
Then
\begin{align*}
x+y=(a_1+b_1)x_1+\cdots+(a_n+b_n)x_n,
\end{align*}
so
\begin{align*}
[x+y]_B=(a_1+b_1,\dots,a_n+b_n)=[x]_B+[y]_B.
\end{align*}
Also, for $c\in F$,
\begin{align*}
cx=(ca_1)x_1+\cdots+(ca_n)x_n,
\end{align*}
hence
\begin{align*}
[cx]_B=(ca_1,\dots,ca_n)=c[x]_B.
\end{align*}
Therefore $\Phi$ is linear.

Next we show injectivity. Suppose $\Phi(x)=\Phi(y)$. Then
\begin{align*}
[x]_B=[y]_B.
\end{align*}
This means the coordinates of $x$ and $y$ with respect to $B$ are identical. Since basis expansions are unique, it follows that $x=y$. Hence $\Phi$ is injective.

Now we show surjectivity. Let $(a_1,\dots,a_n)\in F^n$. Define
\begin{align*}
x=a_1x_1+\cdots+a_nx_n\in V.
\end{align*}
Then by definition,
\begin{align*}
\Phi(x)=[x]_B=(a_1,\dots,a_n).
\end{align*}
Thus every vector in $F^n$ is in the image of $\Phi$, so $\Phi$ is surjective.
Note that in linear algebra, “linear” = “homomorphism.” iT is exactly what it means for $\Phi$ to be a homomorphism of vector spaces.

Since $\Phi$ is linear, injective, and surjective, it is a linear isomorphism.
}}

\bigskip

\clearpage
\hfill \break  $~\!\!$ \dotfill \medskip \\
\textbf{Question.}\\
Let $B, B'$ be two ordered bases of $V$. Show that there exists an invertible matrix $Q$ such that:
\begin{align*}
[x]_{B'} = Q [x]_B
\end{align*}
and identify $Q$ as the change of coordinate matrix.

\noindent\\
\textbf{Solution.}\\
\quad\\
\noindent\fbox{ \parbox{\textwidth}{
Let
\begin{align*}
B=\{x_1,\dots,x_n\},\qquad B'=\{x_1',\dots,x_n'\}.
\end{align*}
Consider the two coordinate isomorphisms
\begin{align*}
\Phi_B:V\to F^n,\qquad \Phi_B(x)=[x]_B,
\end{align*}
and
\begin{align*}
\Phi_{B'}:V\to F^n,\qquad \Phi_{B'}(x)=[x]_{B'}.
\end{align*}
Both are linear isomorphisms. Therefore the composition
\begin{align*}
Q:=\Phi_{B'}\circ \Phi_B^{-1}:F^n\to F^n
\end{align*}
is a linear isomorphism. Hence $Q$ is represented by an invertible $n\times n$ matrix.

Now let $x\in V$. Then
\begin{align*}
Q[x]_B
&=(\Phi_{B'}\circ \Phi_B^{-1})([x]_B) \\
&=\Phi_{B'}(x) \\
&=[x]_{B'}.
\end{align*}
So
\begin{align*}
[x]_{B'}=Q[x]_B.
\end{align*}

Here, this matrix $Q$ is identified as the change of coordinate matrix from $B$-coordinates to $B'$-coordinates.

Equivalently, the $j$th column of $Q$ is the coordinate vector of the basis vector $x_j$ with respect to $B'$, namely
\begin{align*}
Q=\begin{bmatrix}
[x_1]_{B'} & [x_2]_{B'} & \cdots & [x_n]_{B'}
\end{bmatrix}.
\end{align*}
Since $Q$ comes from a linear isomorphism, it is invertible.
}}

\bigskip

\clearpage
\hfill \break  $~\!\!$ \dotfill \medskip \\
\textbf{Question.}\\
Let $T: V \to V$ be linear and let $B, B'$ be bases. Show that:
\begin{align*}
[T]_{B'} = Q [T]_B Q^{-1}
\end{align*}
where $Q$ is the change of coordinate matrix from $B'$ to $B$.

\noindent\\
\textbf{Solution.}\\
\quad\\
\noindent\fbox{ \parbox{\textwidth}{
Let
\begin{align*}
\Phi_B(x)=[x]_B,\qquad \Phi_{B'}(x)=[x]_{B'}.
\end{align*}
Then the matrix of $T$ with respect to $B$ is the unique matrix $[T]_B$ satisfying
\begin{align*}
[T(x)]_B=[T]_B[x]_B\qquad \text{for all }x\in V.
\end{align*}
Similarly,
\begin{align*}
[T(x)]_{B'}=[T]_{B'}[x]_{B'}.
\end{align*}

Let $Q$ be the change of coordinate matrix from $B'$ to $B$, so
\begin{align*}
[x]_B=Q[x]_{B'}.
\end{align*}
Hence
\begin{align*}
[x]_{B'}=Q^{-1}[x]_B.
\end{align*}

Now for any $x\in V$,
\begin{align*}
[T(x)]_{B'}
&=Q^{-1}[T(x)]_B \\
&=Q^{-1}[T]_B[x]_B \\
&=Q^{-1}[T]_BQ[x]_{B'}.
\end{align*}
Since also
\begin{align*}
[T(x)]_{B'}=[T]_{B'}[x]_{B'},
\end{align*}
we conclude that
\begin{align*}
[T]_{B'}=Q^{-1}[T]_BQ.
\end{align*}

If instead one defines $Q$ to be the change of coordinate matrix from $B$ to $B'$, then the formula becomes
\begin{align*}
[T]_{B'}=Q[T]_BQ^{-1}.
\end{align*}
where the exact placement of $Q$ and $Q^{-1}$ depends on the convention for the change of coordinate matrix.
}}

\bigskip

\clearpage
\hfill \break  $~\!\!$ \dotfill \medskip \\
\textbf{Question.}\\
Fill in the missing steps to prove that if $T,S$ are linear, then $T+S$ is linear:
\begin{align*}
(T+S)(x+y) &= (T+S)(x) + (T+S)(y), \\
(T+S)(ax) &= a(T+S)(x).
\end{align*}

\noindent\\
\textbf{Solution.}\\
\quad\\
\noindent\fbox{ \parbox{\textwidth}{
Let $x,y\in X$ and $a\in F$. Then
\begin{align*}
(T+S)(x+y)
&=T(x+y)+S(x+y) \\
&=(T(x)+T(y))+(S(x)+S(y)) \\
&=(T(x)+S(x))+(T(y)+S(y)) \\
&=(T+S)(x)+(T+S)(y).
\end{align*}
Also,
\begin{align*}
(T+S)(ax)
&=T(ax)+S(ax) \\
&=aT(x)+aS(x) \\
&=a(T(x)+S(x)) \\
&=a(T+S)(x).
\end{align*}
Hence $T+S$ is linear.
}}

\bigskip

\clearpage
\hfill \break  $~\!\!$ \dotfill \medskip \\
\textbf{Question.}\\
Show that checking closure under addition and scalar multiplication is sufficient to prove $W \subset V$ is a subspace. Explain why $0 \in W$ automatically follows.

\quad\\
\noindent
\textbf{Solution.}\\
\quad\\
\noindent\fbox{ \parbox{\textwidth}{
Assume $W\neq \varnothing$, and assume:
\begin{align*}
x,y\in W &\implies x+y\in W, \\
a\in F,\ x\in W &\implies ax\in W.
\end{align*}
We show that $W$ is a subspace.

Pick any $w\in W$; this is possible because $W\neq\varnothing$. Since $0\in F$, closure under scalar multiplication gives
\begin{align*}
0\cdot w=0\in W.
\end{align*}
Thus the zero vector belongs to $W$ automatically.

Next, if $w\in W$, then since $-1\in F$,
\begin{align*}
(-1)w=-w\in W.
\end{align*}
So additive inverses are also in $W$.

Because all remaining vector space axioms are inherited from $V$, it follows that $W$ is a subspace of $V$.

Therefore, for a nonempty subset $W$, closure under addition and scalar multiplication is enough to conclude that $W$ is a subspace.
}}

\bigskip

\clearpage
\hfill \break  $~\!\!$ \dotfill \medskip \\
\textbf{Question.}\\
Suppose
\begin{align*}
x = a_1 x_1 + \cdots + a_n x_n = b_1 x_1 + \cdots + b_n x_n.
\end{align*}
Show that $a_i = b_i$ for all $i$ using linear independence.

\noindent\\
\textbf{Solution.}\\
\quad\\
\noindent\fbox{ \parbox{\textwidth}{
Subtract the two expressions:
\begin{align*}
0
&=(a_1x_1+\cdots+a_nx_n)-(b_1x_1+\cdots+b_nx_n) \\
&=(a_1-b_1)x_1+\cdots+(a_n-b_n)x_n.
\end{align*}
If $\{x_1,\dots,x_n\}$ is linearly independent, then the only way a linear combination of these vectors can equal $0$ is if every coefficient is $0$. Hence
\begin{align*}
a_1-b_1=0,\quad \dots,\quad a_n-b_n=0.
\end{align*}
Therefore
\begin{align*}
a_i=b_i\qquad \text{for all }i.
\end{align*}
}}

\bigskip

\clearpage
\hfill \break  $~\!\!$ \dotfill \medskip \\
\textbf{Question.}\\
Prove that the coordinate map $x \mapsto [x]_B$ is linear, injective, and surjective.

\noindent\\
\textbf{Solution.}\\
\quad\\
\noindent\fbox{ \parbox{\textwidth}{
Let $B=\{x_1,\dots,x_n\}$. Define $\Phi(x)=[x]_B$.

To prove linearity, write
\begin{align*}
x=a_1x_1+\cdots+a_nx_n,\qquad y=b_1x_1+\cdots+b_nx_n.
\end{align*}
Then
\begin{align*}
\Phi(x+y)=[x+y]_B=(a_1+b_1,\dots,a_n+b_n)=\Phi(x)+\Phi(y),
\end{align*}
and for $c\in F$,
\begin{align*}
\Phi(cx)=[cx]_B=(ca_1,\dots,ca_n)=c\Phi(x).
\end{align*}
So $\Phi$ is linear.

To prove injectivity, suppose $\Phi(x)=\Phi(y)$. Then $[x]_B=[y]_B$, so $x$ and $y$ have the same coordinates with respect to $B$. By uniqueness of basis expansions, $x=y$.

To prove surjectivity, let $(a_1,\dots,a_n)\in F^n$. Set
\begin{align*}
x=a_1x_1+\cdots+a_nx_n.
\end{align*}
Then
\begin{align*}
\Phi(x)=[x]_B=(a_1,\dots,a_n).
\end{align*}
Hence every vector in $F^n$ is hit by $\Phi$.

Therefore $\Phi$ is linear, injective, and surjective.
}}

\bigskip


\clearpage
\hfill \break  $~\!\!$ \dotfill \medskip \\
\textbf{Question.}\\
Construct the matrix $Q$ explicitly by expressing each vector of $B'$ in terms of $B$. Show that $Q$ is invertible.

\noindent\\
\textbf{Solution.}\\
\quad\\
\noindent\fbox{ \parbox{\textwidth}{
Let
\begin{align*}
B=\{x_1,\dots,x_n\},\qquad B'=\{x_1',\dots,x_n'\}.
\end{align*}
Because $B$ is a basis, each $x_j'$ can be written uniquely as
\begin{align*}
x_j'=a_{1j}x_1+\cdots+a_{nj}x_n.
\end{align*}
Define the matrix
\begin{align*}
Q=\begin{bmatrix}
a_{11} & a_{12} & \cdots & a_{1n} \\
a_{21} & a_{22} & \cdots & a_{2n} \\
\vdots & \vdots & \ddots & \vdots \\
a_{n1} & a_{n2} & \cdots & a_{nn}
\end{bmatrix}.
\end{align*}
The $j$th column of $Q$ is exactly $[x_j']_B$.

Now let
\begin{align*}
x=c_1x_1'+\cdots+c_nx_n'.
\end{align*}
Substituting the expressions for the $x_j'$ in terms of $B$ shows that the coordinate vector $[x]_B$ is obtained from $[x]_{B'}$ by multiplication by $Q$:
\begin{align*}
[x]_B=Q[x]_{B'}.
\end{align*}
So $Q$ is the change of coordinate matrix from $B'$-coordinates to $B$-coordinates.

To show $Q$ is invertible, suppose $Qc=0$ for some $c=(c_1,\dots,c_n)^T$. Then
\begin{align*}
0=Qc=[c_1x_1'+\cdots+c_nx_n']_B.
\end{align*}
Hence
\begin{align*}
c_1x_1'+\cdots+c_nx_n'=0.
\end{align*}
Since $B'$ is linearly independent, all $c_j=0$. Thus $\ker(Q)=\{0\}$, so $Q$ is injective. Because $Q$ is an $n\times n$ matrix, injective implies invertible. Therefore $Q$ is invertible.
}}

\bigskip


\clearpage
\hfill \break  $~\!\!$ \dotfill \medskip \\
\textbf{Question.}\\
Show that the matrix representation $[T]_B$ satisfies
\begin{align*}
[T(x)]_B = [T]_B [x]_B
\end{align*}
for all $x \in V$.

\noindent\\
\textbf{Solution.}\\
\quad\\
\noindent\fbox{ \parbox{\textwidth}{
Let $B=\{x_1,\dots,x_n\}$ be a basis of $V$. By definition, the $j$th column of $[T]_B$ is the coordinate vector of $T(x_j)$ with respect to $B$:
\begin{align*}
[T(x_j)]_B=\text{column }j\text{ of }[T]_B.
\end{align*}

Now let
\begin{align*}
x=a_1x_1+\cdots+a_nx_n.
\end{align*}
Then
\begin{align*}
[x]_B=
\begin{bmatrix}
a_1\\
\vdots\\
a_n
\end{bmatrix}.
\end{align*}
Since $T$ is linear,
\begin{align*}
T(x)=a_1T(x_1)+\cdots+a_nT(x_n).
\end{align*}
Taking coordinates with respect to $B$ gives
\begin{align*}
[T(x)]_B
&=a_1[T(x_1)]_B+\cdots+a_n[T(x_n)]_B.
\end{align*}
But multiplying the matrix $[T]_B$ by the vector $[x]_B$ produces exactly this same linear combination of the columns of $[T]_B$. Hence
\begin{align*}
[T(x)]_B=[T]_B[x]_B.
\end{align*}
This is the defining relation between a linear map and its matrix representation with respect to the basis $B$.
}}

\bigskip

\clearpage
\hfill \break  $~\!\!$ \dotfill \medskip \\
\textbf{Question.}\\
Show that $W \neq \varnothing$ and the single condition
\[
c \in F,\; x, y \in W \implies cx + y \in W
\]
is sufficient to prove $W$ is a subspace of $V$.

\noindent\\
\textbf{Solution.}\\
\quad\\
\noindent\fbox{ \parbox{\textwidth}{
\textbf{Zero vector:} Since $W \neq \varnothing$, pick any $w \in W$. Take $c = 0$, $x = y = w$:
\[
0 \cdot w + w = w \in W.
\]
Now take $c = -1$, $x = y = w$:
\[
(-1)w + w = 0 \in W.
\]

\textbf{Closure under addition:} Let $x, y \in W$. Take $c = 1$:
\[
1 \cdot x + y = x + y \in W.
\]

\textbf{Closure under scalar multiplication:} Let $x \in W$, $a \in F$. Since $0 \in W$, take $y = 0$, $c = a$:
\[
a \cdot x + 0 = ax \in W.
\]

Hence all three subspace conditions are satisfied, so $W$ is a subspace.

\hfill $\square$
}}

\bigskip

\clearpage
\hfill \break  $~\!\!$ \dotfill \medskip \\
\textbf{Question.}\\
Let $W_1, W_2$ be subspaces of a vector space $V$ over $F$. Show that $W_1 \cap W_2$ is a subspace of $V$.

\noindent\\
\textbf{Solution.}\\
\quad\\
\noindent\fbox{ \parbox{\textwidth}{
\textbf{Zero vector:} Since $W_1$ and $W_2$ are subspaces, $0 \in W_1$ and $0 \in W_2$, so $0 \in W_1 \cap W_2$.

\textbf{Closure under addition:} Let $x, y \in W_1 \cap W_2$. Then $x, y \in W_1$ and $x, y \in W_2$. Since $W_1$ is a subspace, $x + y \in W_1$. Since $W_2$ is a subspace, $x + y \in W_2$. Hence $x + y \in W_1 \cap W_2$.

\textbf{Closure under scalar multiplication:} Let $a \in F$, $x \in W_1 \cap W_2$. Then $x \in W_1 \implies ax \in W_1$ and $x \in W_2 \implies ax \in W_2$. Hence $ax \in W_1 \cap W_2$.

Therefore $W_1 \cap W_2$ is a subspace of $V$.

\hfill $\square$
}}

\bigskip

\clearpage
\hfill \break  $~\!\!$ \dotfill \medskip \\
\textbf{Question.}\\
Let $W_1, W_2$ be subspaces of $V$. Recall $W_1 + W_2 = \{x_1 + x_2 : x_1 \in W_1,\; x_2 \in W_2\}$. Show:
\begin{enumerate}
    \item $W_1 + W_2$ is a subspace of $V$ that contains both $W_1$ and $W_2$.
    \item Any subspace of $V$ that contains both $W_1$ and $W_2$ must also contain $W_1 + W_2$.
\end{enumerate}
Conclude that $W_1 + W_2$ is the smallest subspace containing both $W_1$ and $W_2$.

\noindent\\
\textbf{Solution.}\\
\quad\\
\noindent\fbox{ \parbox{\textwidth}{
\textbf{Part 1:}

\textit{$W_1 + W_2$ is a subspace:}

\textbf{Zero:} $0 = 0 + 0$ with $0 \in W_1$ and $0 \in W_2$, so $0 \in W_1 + W_2$.

\textbf{Addition:} Let $u, v \in W_1 + W_2$. Write $u = x_1 + x_2$ and $v = y_1 + y_2$ with $x_1, y_1 \in W_1$ and $x_2, y_2 \in W_2$. Then $u + v = (x_1 + y_1) + (x_2 + y_2)$ with $x_1 + y_1 \in W_1$ and $x_2 + y_2 \in W_2$, so $u + v \in W_1 + W_2$.

\textbf{Scalar multiplication:} Let $a \in F$, $u = x_1 + x_2 \in W_1 + W_2$. Then $au = ax_1 + ax_2$ with $ax_1 \in W_1$ and $ax_2 \in W_2$, so $au \in W_1 + W_2$.

\textit{Contains both:} For any $w_1 \in W_1$, $w_1 = w_1 + 0 \in W_1 + W_2$. Similarly $W_2 \subseteq W_1 + W_2$.

\textbf{Part 2:} Let $U$ be any subspace containing both $W_1$ and $W_2$. Let $u \in W_1 + W_2$, so $u = x_1 + x_2$ with $x_1 \in W_1 \subseteq U$ and $x_2 \in W_2 \subseteq U$. Since $U$ is a subspace, $x_1 + x_2 \in U$. Hence $W_1 + W_2 \subseteq U$.

\textbf{Conclusion:} $W_1 + W_2$ contains both $W_1$ and $W_2$, and is contained in every subspace that does so. Therefore it is the smallest such subspace.

\hfill $\square$
}}

\bigskip
\begin{center}
\rule{0.5\textwidth}{1pt}\\[6pt]
\textbf{Homework from Lecture on February 10, 2026}\\[2pt]
\rule{0.5\textwidth}{1pt}
\end{center}
\bigskip

\clearpage
\hfill \break  $~\!\!$ \dotfill \medskip \\
\textbf{Question.}\\
Show that $\mathbb{R} \cong \{x + 0i : x \in \mathbb{R}\} \subset \mathbb{C}$ is a subfield of $\mathbb{C}$.

\noindent\\
\textbf{Solution.}\\
\quad\\
\noindent\fbox{ \parbox{\textwidth}{
We already proved $\mathbb{C}$ is a field. We identify $\mathbb{R}$ with $\{x + 0i : x \in \mathbb{R}\} \subset \mathbb{C}$.

\textbf{Closed under addition:} If $a, b \in \mathbb{R}$, then $(a + 0i) + (b + 0i) = (a+b) + 0i \in \mathbb{R}$.

\textbf{Closed under multiplication:} $(a + 0i)(b + 0i) = ab + 0i \in \mathbb{R}$.

\textbf{Contains $0$ and $1$:} $0 = 0 + 0i$ and $1 = 1 + 0i$ are in $\mathbb{R}$.

\textbf{Additive inverses:} $-(a + 0i) = (-a) + 0i \in \mathbb{R}$.

\textbf{Multiplicative inverses:} For $a \neq 0$, $(a + 0i)^{-1} = (1/a) + 0i \in \mathbb{R}$.

Hence $\mathbb{R}$ is a subfield of $\mathbb{C}$: it is a subset that is itself a field under the same operations.

\hfill $\square$
}}

\bigskip

\clearpage
\hfill \break  $~\!\!$ \dotfill \medskip \\
\textbf{Question.}\\
Show that $|\bar{z}| = |z|$ and $z\bar{z} = |z|^2$ for all $z \in \mathbb{C}$.

\noindent\\
\textbf{Solution.}\\
\quad\\
\noindent\fbox{ \parbox{\textwidth}{
Let $z = x + iy$ with $x, y \in \mathbb{R}$. Then $\bar{z} = x - iy$.

\textbf{$|\bar{z}| = |z|$:}
\[
|\bar{z}| = |x - iy| = \sqrt{x^2 + (-y)^2} = \sqrt{x^2 + y^2} = |z|.
\]

\textbf{$z\bar{z} = |z|^2$:}
\begin{align*}
z\bar{z} &= (x + iy)(x - iy) = x^2 - (iy)^2 = x^2 - i^2 y^2 = x^2 + y^2.
\end{align*}
Since $|z|^2 = (\sqrt{x^2+y^2})^2 = x^2 + y^2$, we have $z\bar{z} = |z|^2$.

\hfill $\square$
}}

\bigskip

\clearpage
\hfill \break  $~\!\!$ \dotfill \medskip \\
\textbf{Question.}\\
Show that $d(z_1, z_2) = |z_1 - z_2|$ defines a metric on $\mathbb{C}$. That is, verify non-negativity, symmetry, and the triangle inequality.

\noindent\\
\textbf{Solution.}\\
\quad\\
\noindent\fbox{ \parbox{\textwidth}{
\textbf{Non-negativity:} $d(z_1, z_2) = |z_1 - z_2| \ge 0$ since the modulus is non-negative. Also $d(z_1, z_2) = 0 \iff |z_1 - z_2| = 0 \iff z_1 - z_2 = 0 \iff z_1 = z_2$.

\textbf{Symmetry:}
\[
d(z_1, z_2) = |z_1 - z_2| = |{-(z_2 - z_1)}| = |z_2 - z_1| = d(z_2, z_1).
\]

\textbf{Triangle inequality:} For $z_1, z_2, z_3 \in \mathbb{C}$, write $z_1 - z_3 = (z_1 - z_2) + (z_2 - z_3)$. By the triangle inequality for $|\cdot|$:
\[
d(z_1, z_3) = |z_1 - z_3| = |(z_1 - z_2) + (z_2 - z_3)| \le |z_1 - z_2| + |z_2 - z_3| = d(z_1, z_2) + d(z_2, z_3).
\]

Hence $d$ is a metric on $\mathbb{C}$.

\hfill $\square$
}}

\bigskip

\clearpage
\hfill \break  $~\!\!$ \dotfill \medskip \\
\textbf{Question.}\\
Show by induction that for all $z_1, \dots, z_n \in \mathbb{C}$:
\[
|z_1 + z_2 + \cdots + z_n| \le |z_1| + |z_2| + \cdots + |z_n|.
\]

\noindent\\
\textbf{Solution.}\\
\quad\\
\noindent\fbox{ \parbox{\textwidth}{
\textbf{Base case ($n = 2$):} This is the triangle inequality for $|\cdot|$ on $\mathbb{C}$:
\[
|z_1 + z_2| \le |z_1| + |z_2|.
\]

\textbf{Inductive step:} Assume the statement holds for some $n \ge 2$:
\[
|z_1 + \cdots + z_n| \le |z_1| + \cdots + |z_n|.
\]
For $n+1$:
\begin{align*}
|z_1 + \cdots + z_n + z_{n+1}|
&= |(z_1 + \cdots + z_n) + z_{n+1}| \\
&\le |z_1 + \cdots + z_n| + |z_{n+1}| \quad \text{(triangle inequality for $n = 2$)} \\
&\le (|z_1| + \cdots + |z_n|) + |z_{n+1}| \quad \text{(inductive hypothesis)} \\
&= |z_1| + \cdots + |z_{n+1}|.
\end{align*}

By induction, the result holds for all $n \ge 2$.

\hfill $\square$
}}

\bigskip
\begin{center}
\rule{0.5\textwidth}{1pt}\\[6pt]
\textbf{Homework from Lecture on March 3, 2026}\\[2pt]
\rule{0.5\textwidth}{1pt}
\end{center}
\bigskip

\clearpage
\hfill \break  $~\!\!$ \dotfill \medskip \\
\textbf{Question.}\\
Let $T, S : X \to Y$ be linear operators. Show that $\alpha T + \beta S$ is a linear operator for any scalars $\alpha, \beta$, and that the collection of all linear operators from $X$ to $Y$ forms a vector space.

\vspace{0.5cm}
\noindent
\textbf{Solution.}\\
\quad\\
\noindent\fbox{ \parbox{\textwidth}{
Let $u,v \in X$ and let $c \in F$. Define $(\alpha T+\beta S):X\to Y$ by
\begin{align*}
(\alpha T+\beta S)(x)=\alpha T(x)+\beta S(x).
\end{align*}
This is the pointwise definition of addition and scalar multiplication of functions. We can do this because $T(x)$ and $S(x)$ are elements of $Y$ since $T,S:X\to Y$; $Y$ is a vector space over $F$, so $\alpha T(x)$ and $\beta S(x)$ are in $Y$; and since $Y$ is closed under addition, $\alpha T(x)+\beta S(x)\in Y$. \\

We check linearity:
\begin{align*}
(\alpha T+\beta S)(u+v)
&=\alpha T(u+v)+\beta S(u+v) \\
&=\alpha(T(u)+T(v))+\beta(S(u)+S(v)) \\
&=(\alpha T(u)+\beta S(u))+(\alpha T(v)+\beta S(v)) \\
&=(\alpha T+\beta S)(u)+(\alpha T+\beta S)(v),
\end{align*}
and
\begin{align*}
(\alpha T+\beta S)(cu)
&=\alpha T(cu)+\beta S(cu) \\
&=\alpha cT(u)+\beta cS(u) \\
&=c(\alpha T(u)+\beta S(u)) \\
&=c(\alpha T+\beta S)(u).
\end{align*}
Hence $\alpha T+\beta S$ is linear.

Now let $\mathcal{L}(X,Y)$ denote the set of all linear operators from $X$ to $Y$. We show it is a vector space over $F$ with the usual addition and scalar multiplication.

First, by the calculation above, $\mathcal{L}(X,Y)$ is closed under addition and scalar multiplication. Also the zero map $0:X\to Y$, defined by $0(x)=0$ for all $x\in X$, is linear, so $0\in \mathcal{L}(X,Y)$. If $T\in \mathcal{L}(X,Y)$, then $(-1)T=-T$ is linear, so additive inverses exist.

The remaining vector space axioms follow pointwise from the corresponding axioms in $Y$. For example, for $T,S,R\in \mathcal{L}(X,Y)$ and $x\in X$,
\begin{align*}
((T+S)+R)(x)=T(x)+S(x)+R(x)=(T+(S+R))(x),
\end{align*}
so $(T+S)+R=T+(S+R)$. Similarly $T+S=S+T$, and for $a,b\in F$,
\begin{align*}
(a+b)T=aT+bT,\qquad a(T+S)=aT+aS,\qquad (ab)T=a(bT),\qquad 1T=T.
\end{align*}
Therefore $\mathcal{L}(X,Y)$ is a vector space.
}}

\bigskip

\clearpage
\hfill \break  $~\!\!$ \dotfill \medskip \\
\textbf{Question.}\\
Let $(X, d)$ be a metric space. Carefully define what it means for a sequence $\{x_n\}$ in $X$ to be \emph{convergent} and what it means for it to be \emph{Cauchy}. Then prove that every convergent sequence in any metric space is Cauchy. (Use $d$, not $\|\cdot\|$.)

\noindent\\
\textbf{Solution.}\\
\quad\\
\noindent\fbox{ \parbox{\textwidth}{
\textbf{Definitions:}

A sequence $\{x_n\}$ in $(X, d)$ \textbf{converges} to $x \in X$ if for every $\varepsilon > 0$ there exists $N \in \mathbb{N}$ such that
\[
n > N \implies d(x_n, x) < \varepsilon.
\]

A sequence $\{x_n\}$ in $(X, d)$ is \textbf{Cauchy} if for every $\varepsilon > 0$ there exists $N \in \mathbb{N}$ such that
\[
n, m > N \implies d(x_n, x_m) < \varepsilon.
\]

\textbf{Proof that convergent implies Cauchy:}

Suppose $x_n \to x$. Let $\varepsilon > 0$. By convergence, there exists $N$ such that $d(x_n, x) < \varepsilon/2$ for all $n > N$.

Then for $n, m > N$, the triangle inequality gives
\begin{align*}
d(x_n, x_m) \le d(x_n, x) + d(x, x_m) < \frac{\varepsilon}{2} + \frac{\varepsilon}{2} = \varepsilon.
\end{align*}
Hence $\{x_n\}$ is Cauchy.

Note: the converse (Cauchy $\Rightarrow$ convergent) holds only in \emph{complete} metric spaces. That is exactly the definition of completeness.

\hfill $\square$
}}

\bigskip
\begin{center}
\rule{0.5\textwidth}{1pt}\\[6pt]
\textbf{Homework from Lecture on March 10, 2026}\\[2pt]
\rule{0.5\textwidth}{1pt}
\end{center}
\bigskip

\clearpage
\hfill \break  $~\!\!$ \dotfill \medskip \\
\textbf{Question.}\\
Let $X = (X, \langle \cdot, \cdot \rangle)$ be an inner product space over $\mathbb{R}$. The Pythagorean theorem states: if $x \perp y$, then $\|x+y\|^2 = \|x\|^2 + \|y\|^2$. Using the identity
\[
\langle x,y \rangle = \tfrac{1}{2}\bigl(\|x+y\|^2 - \|x\|^2 - \|y\|^2\bigr),
\]
show that this is the law of cosines in $X = \mathbb{R}^2$ with the standard dot product.
\noindent\\

\noindent\\
\textbf{Solution.}\\
\quad\\
\noindent\fbox{ \parbox{\textwidth}{
In $\mathbb{R}^2$ with the standard inner product $\langle x, y \rangle = x \cdot y = \|x\|\|y\|\cos\theta$, the identity
\[
\|x+y\|^2 = \|x\|^2 + \|y\|^2 + 2\langle x, y \rangle
\]
becomes
\[
\|x+y\|^2 = \|x\|^2 + \|y\|^2 + 2\|x\|\|y\|\cos\theta.
\]
Let $a = \|x\|$, $b = \|y\|$, and $c = \|x+y\|$. Then
\[
c^2 = a^2 + b^2 + 2ab\cos\theta.
\]
Replacing $y$ by $-y$ (so the angle between $x$ and $-y$ is $\pi - \theta$), the side opposite $\theta$ has length $\|x - y\|$ and
\[
\|x - y\|^2 = a^2 + b^2 - 2ab\cos\theta,
\]
which is the classical law of cosines. When $\theta = \pi/2$ (i.e., $x \perp y$), we recover the Pythagorean theorem: $c^2 = a^2 + b^2$.

\hfill $\square$
}}

\bigskip

\clearpage
\hfill \break  $~\!\!$ \dotfill \medskip \\
\textbf{Question.}\\
Let $X = (X, \langle \cdot, \cdot \rangle)$ be an inner product space over $\mathbb{R}$. For non-zero $x, y \in X$, define
\[
\cos\theta = \frac{\langle x, y \rangle}{\|x\|\|y\|}.
\]
Show that when $X = \mathbb{R}^2$ with the dot product, $\theta$ is the usual angle between $x$ and $y$.

\noindent\\
\textbf{Solution.}\\
\quad\\
\noindent\fbox{ \parbox{\textwidth}{
By the Cauchy--Schwarz inequality, $|\langle x,y\rangle| \le \|x\|\|y\|$, so
\[
-1 \le \frac{\langle x, y \rangle}{\|x\|\|y\|} \le 1.
\]
Since $\cos$ is one-to-one on $[0,\pi]$, there exists a unique $\theta \in [0,\pi]$ such that $\cos\theta = \langle x,y\rangle / (\|x\|\|y\|)$.

In $\mathbb{R}^2$, let $x = (x_1, x_2)$ and $y = (y_1, y_2)$. Write $x = \|x\|(\cos\alpha, \sin\alpha)$ and $y = \|y\|(\cos\beta, \sin\beta)$ where $\alpha, \beta$ are the angles these vectors make with the positive $x$-axis. Then
\begin{align*}
\langle x, y \rangle &= \|x\|\|y\|(\cos\alpha\cos\beta + \sin\alpha\sin\beta) \\
&= \|x\|\|y\|\cos(\alpha - \beta).
\end{align*}
So $\langle x,y\rangle / (\|x\|\|y\|) = \cos(\alpha - \beta)$, and $\theta = |\alpha - \beta|$ (up to the range $[0,\pi]$), which is exactly the geometric angle between $x$ and $y$.

\hfill $\square$
}}

\bigskip

\clearpage
\hfill \break  $~\!\!$ \dotfill \medskip \\
\textbf{Question.}\\
(\textbf{Polarization identity over $\mathbb{R}$.}) Let $X = (X, \langle \cdot, \cdot \rangle)$ be an inner product space over $\mathbb{R}$. Show that for all $x, y \in X$:
\[
\langle x, y \rangle = \frac{1}{4}\bigl(\|x+y\|^2 - \|x-y\|^2\bigr).
\]

\noindent\\
\textbf{Solution.}\\
\quad\\
\noindent\fbox{ \parbox{\textwidth}{
By the law of cosines (Proposition from lecture):
\begin{align*}
\|x+y\|^2 &= \|x\|^2 + \|y\|^2 + 2\,\text{Re}\,\langle x, y \rangle, \\
\|x-y\|^2 &= \|x\|^2 + \|y\|^2 - 2\,\text{Re}\,\langle x, y \rangle.
\end{align*}
Since the field is $\mathbb{R}$, $\text{Re}\,\langle x, y \rangle = \langle x, y \rangle$. Subtracting:
\begin{align*}
\|x+y\|^2 - \|x-y\|^2 &= 4\langle x, y \rangle.
\end{align*}
Dividing by $4$:
\[
\langle x, y \rangle = \frac{1}{4}\bigl(\|x+y\|^2 - \|x-y\|^2\bigr).
\]

\hfill $\square$
}}

\bigskip

\clearpage
\hfill \break  $~\!\!$ \dotfill \medskip \\
\textbf{Question.}\\
Let $V, W$ be inner product spaces over $\mathbb{C}$. A linear map $T: V \to W$ is an isometry if $\langle T(u), T(v) \rangle_W = \langle u, v \rangle_V$ for all $u, v$. Show that $T$ is an isometry if and only if $\|T(v)\|_W = \|v\|_V$ for all $v \in V$.
(\textit{Hint: use the polarization identity for $\mathbb{C}$.})

\noindent\\
\textbf{Solution.}\\
\quad\\
\noindent\fbox{ \parbox{\textwidth}{
$(\Rightarrow)$: If $\langle T(u), T(v) \rangle = \langle u, v \rangle$ for all $u, v$, then taking $u = v$ gives $\|T(v)\|^2 = \langle T(v), T(v) \rangle = \langle v, v \rangle = \|v\|^2$, so $\|T(v)\| = \|v\|$.

$(\Leftarrow)$: Suppose $\|T(v)\| = \|v\|$ for all $v$. We use the polarization identity over $\mathbb{C}$:
\[
\langle u, v \rangle = \frac{1}{4}\bigl(\|u+v\|^2 - \|u-v\|^2 + i\|u+iv\|^2 - i\|u-iv\|^2\bigr).
\]
Since $T$ is linear:
\begin{align*}
\langle T(u), T(v) \rangle &= \frac{1}{4}\bigl(\|T(u)+T(v)\|^2 - \|T(u)-T(v)\|^2 \\
&\qquad + i\|T(u)+iT(v)\|^2 - i\|T(u)-iT(v)\|^2\bigr) \\
&= \frac{1}{4}\bigl(\|T(u+v)\|^2 - \|T(u-v)\|^2 \\
&\qquad + i\|T(u+iv)\|^2 - i\|T(u-iv)\|^2\bigr).
\end{align*}
By the norm-preserving hypothesis, $\|T(w)\|^2 = \|w\|^2$ for all $w$. Substituting:
\[
= \frac{1}{4}\bigl(\|u+v\|^2 - \|u-v\|^2 + i\|u+iv\|^2 - i\|u-iv\|^2\bigr) = \langle u, v \rangle.
\]

\hfill $\square$
}}

\bigskip

\clearpage
\hfill \break  $~\!\!$ \dotfill \medskip \\
\textbf{Question.}\\
Let $(X, d_X)$ and $(Y, d_Y)$ be metric spaces. Show that if $f: X \to Y$ is Lipschitz continuous, then $f$ is continuous.

\noindent\\
\textbf{Solution.}\\
\quad\\
\noindent\fbox{ \parbox{\textwidth}{
Suppose $f$ is Lipschitz: there exists $L \ge 0$ such that
\[
d_Y(f(x), f(y)) \le L\, d_X(x, y) \quad \text{for all } x, y \in X.
\]
We show $f$ is continuous at every $x_0 \in X$. Let $\varepsilon > 0$. If $L = 0$, then $f$ is constant and continuity is trivial. If $L > 0$, set $\delta = \varepsilon / L$. Then for any $x$ with $d_X(x, x_0) < \delta$:
\[
d_Y(f(x), f(x_0)) \le L\, d_X(x, x_0) < L \cdot \frac{\varepsilon}{L} = \varepsilon.
\]
Hence $f$ is continuous at $x_0$. Since $x_0$ was arbitrary, $f$ is continuous on $X$.

Note: this uses the metric $d$, not the norm $\|\cdot\|$, so it holds in any metric space, not just normed spaces.

\hfill $\square$
}}

\bigskip
\begin{center}
\rule{0.5\textwidth}{1pt}\\[6pt]
\textbf{Homework from Lecture on March 18, 2026}\\[2pt]
\rule{0.5\textwidth}{1pt}
\end{center}
\bigskip

\clearpage
\hfill \break  $~\!\!$ \dotfill \medskip \\
\textbf{Question.}\\
Let $X := (X, \|\cdot \|_X$) and $Y= (Y, \|\cdot \|_Y$) be normed spaces, and let $T: X\rightarrow Y$ be a linear operator. Prove that if $dim(X)<\infty$, then $T$ is bounded.
\noindent\\
\textbf{Solution.}\\
\quad\\
\noindent\fbox{ \parbox{\textwidth}{
Suppose $\dim(X) = n < \infty$, and let $\{e_1, \dots, e_n\}$ be a basis for $X$. Then for any $x \in X$, there exist unique scalars $a_1, \dots, a_n$ such that
\[
x = \sum_{i=1}^n a_i e_i.
\]
By linearity of $T$, we have
\[
T(x) = \sum_{i=1}^n a_i T(e_i).
\]
Using the triangle inequality,
\[
\|T(x)\|_Y \le \sum_{i=1}^n |a_i| \, \|T(e_i)\|_Y.
\]
Let
\[
M := \max_{1 \le i \le n} \|T(e_i)\|_Y.
\]
Then
\[
\|T(x)\|_Y \le M \sum_{i=1}^n |a_i|.
\]

Define a new norm on $X$ based on the basis $\{e_1, \dots, e_n\}$ by
\[
\|x\|_* := \sum_{i=1}^n |a_i|.
\]
Since $X$ is finite-dimensional, all norms on $X$ are equivalent. Hence there exists a constant $C > 0$ such that
\[
\|x\|_* \le C \|x\|_X \quad \text{for all } x \in X.
\]
Therefore,
\[
\|T(x)\|_Y \le M \|x\|_* \le MC \|x\|_X.
\]
Let $K := MC$. Then for all $x \in X$,
\[
\|T(x)\|_Y \le K \|x\|_X.
\]

Thus $T$ is bounded.

\hfill $\square$

}}

\bigskip

\clearpage
\hfill \break  $~\!\!$ \dotfill \medskip \\
\textbf{Question.}\\
Let $T: X\rightarrow Y$ be a linear operator between normed linear spaces. Suppose that $T$ is continuous. Then, show that $T(0)=0$.
\noindent\\
\textbf{Solution.}\\
\quad\\
\noindent\fbox{ \parbox{\textwidth}{
Since $T$ is linear, for any $x \in X$ we have
\[
T(0) = T(x - x) = T(x) - T(x) = 0.
\]
Thus $T(0) = 0$.

\medskip

Alternatively, using continuity: since $T$ is continuous at $0$, for any sequence $\{x_n\}$ in $X$ with $x_n \to 0$, we have $T(x_n) \to T(0)$. Taking $x_n = 0$ for all $n$, we get $T(x_n) = 0$ for all $n$, hence $T(0) = 0$.

\hfill $\square$
}}

\bigskip

\clearpage
\hfill \break  $~\!\!$ \dotfill \medskip \\
\textbf{Question.}\\
Let $T: X \to Y$ be a linear operator between normed linear spaces. Suppose that $T$ is continuous. Then, for $\varepsilon = 1$, there exists $\delta > 0$ such that
\[
\|u - 0\| < \delta \implies \|T(u) - T(0)\| < 1.
\]
Since $T(0) = 0$, this simplifies to
\[
\|u\| < \delta \implies \|T(u)\| < 1.
\]

Show that this implies
\[
\|u\| < \delta \implies \|T(u)\| \le 1.
\]
\noindent\\
\textbf{Solution.}\\
\quad\\
\noindent\fbox{ \parbox{\textwidth}{
Let $u \in X$ with $\|u\| < \delta$. For any $0 < t < 1$, we have
\[
\|t u\| = t \|u\| < \|u\| < \delta.
\]
Hence, by the previous implication,
\[
\|T(tu)\| < 1.
\]
Using linearity of $T$, we obtain
\[
\|t T(u)\| < 1 \quad \Rightarrow \quad t \|T(u)\| < 1.
\]
Thus,
\[
\|T(u)\| < \frac{1}{t}.
\]
Since this holds for all $0 < t < 1$, letting $t \to 1^{-}$ gives
\[
\|T(u)\| \le 1.
\]

\hfill $\square$

}}

\bigskip
\begin{center}
\rule{0.5\textwidth}{1pt}\\[6pt]
\textbf{Homework from Lecture on March 24, 2026}\\[2pt]
\rule{0.5\textwidth}{1pt}
\end{center}
\bigskip

\clearpage
\hfill \break  $~\!\!$ \dotfill \medskip \\
\textbf{Question.}\\
Let $V$ be an inner product space and let $X \subseteq V$ be any subset (not necessarily a subspace). Show that the orthogonal complement
\[
X^\perp = \{ v \in V \mid \langle v, x \rangle = 0 \text{ for all } x \in X \}
\]
is a subspace of $V$.

\noindent\\
\textbf{Solution.}\\
\quad\\
\noindent\fbox{ \parbox{\textwidth}{
We verify the subspace conditions: closure under addition, closure under scalar multiplication, and containment of the zero vector.

\textbf{Zero vector:} For all $x \in X$,
\begin{align*}
\langle 0, x \rangle = 0.
\end{align*}
Hence $0 \in X^\perp$, so $X^\perp$ is non-empty.

\textbf{Closure under addition:} Let $u, v \in X^\perp$. Then for all $x \in X$,
\begin{align*}
\langle u + v, x \rangle &= \langle u, x \rangle + \langle v, x \rangle \\
&= 0 + 0 = 0.
\end{align*}
Here we used linearity of the inner product in the first variable. Hence $u + v \in X^\perp$.

\textbf{Closure under scalar multiplication:} Let $u \in X^\perp$ and $a \in F$. Then for all $x \in X$,
\begin{align*}
\langle a u, x \rangle &= a \langle u, x \rangle \\
&= a \cdot 0 = 0.
\end{align*}
Hence $a u \in X^\perp$.

Since $X^\perp$ contains $0$ and is closed under addition and scalar multiplication, $X^\perp$ is a subspace of $V$.

\hfill $\square$
}}

\bigskip

\clearpage
\hfill \break  $~\!\!$ \dotfill \medskip \\
\textbf{Question.}\\
Let $V$ be an inner product space and let $S \subseteq V$ be a subspace. Show that
\[
S \cap S^\perp = \{0\}.
\]

\noindent\\
\textbf{Solution.}\\
\quad\\
\noindent\fbox{ \parbox{\textwidth}{
First, $0 \in S \cap S^\perp$ because $0$ belongs to every subspace. So $\{0\} \subseteq S \cap S^\perp$.

Now suppose $v \in S \cap S^\perp$. Then:
\begin{itemize}
    \item $v \in S$, and
    \item $v \in S^\perp$, which means $\langle v, s \rangle = 0$ for all $s \in S$.
\end{itemize}

Since $v \in S$, we may take $s = v$ in the second condition:
\begin{align*}
\langle v, v \rangle = 0.
\end{align*}

By the non-degeneracy (positive definiteness) of the inner product, $\langle v, v \rangle = 0$ implies $v = 0$.

Therefore $S \cap S^\perp \subseteq \{0\}$, and hence $S \cap S^\perp = \{0\}$.

\hfill $\square$
}}

\bigskip

\clearpage
\hfill \break  $~\!\!$ \dotfill \medskip \\
\textbf{Question.}\\
Prove the Gram--Schmidt process by induction. Specifically, let $B = (v_1, v_2, \dots)$ be a sequence of non-zero vectors in an inner product space $V$. Define
\begin{align*}
u_1 &= v_1, \\
u_k &= v_k - \sum_{i=1}^{k-1} \frac{\langle v_k, u_i \rangle}{\langle u_i, u_i \rangle}\, u_i, \quad k \ge 2.
\end{align*}
Show by induction on $k$ that for all $k \ge 1$:
\begin{enumerate}
    \item $\{u_1, \dots, u_k\}$ is an orthogonal set of non-zero vectors, and
    \item $\operatorname{span}\{u_1, \dots, u_k\} = \operatorname{span}\{v_1, \dots, v_k\}$.
\end{enumerate}

\noindent\\
\textbf{Solution.}\\
\quad\\
\noindent\fbox{ \parbox{\textwidth}{
\textbf{Base case ($k=1$):}
$u_1 = v_1$. The set $\{u_1\}$ is trivially orthogonal (a single vector is orthogonal to itself vacuously for the pairwise condition). Since $v_1 \neq 0$ by hypothesis, $u_1 \neq 0$. Also, $\operatorname{span}\{u_1\} = \operatorname{span}\{v_1\}$. $\checkmark$

\textbf{Inductive step:} Assume the statement holds for some $k \ge 1$. That is:
\begin{itemize}
    \item $\{u_1, \dots, u_k\}$ is an orthogonal set of non-zero vectors, and
    \item $\operatorname{span}\{u_1, \dots, u_k\} = \operatorname{span}\{v_1, \dots, v_k\}$.
\end{itemize}

We must show the statement holds for $k+1$.

Define
\begin{align*}
u_{k+1} = v_{k+1} - \sum_{i=1}^{k} \frac{\langle v_{k+1}, u_i \rangle}{\langle u_i, u_i \rangle}\, u_i.
\end{align*}

\textbf{Orthogonality:} For each $j \in \{1, \dots, k\}$,
\begin{align*}
\langle u_{k+1}, u_j \rangle
&= \left\langle v_{k+1} - \sum_{i=1}^{k} \frac{\langle v_{k+1}, u_i \rangle}{\langle u_i, u_i \rangle}\, u_i,\; u_j \right\rangle \\
&= \langle v_{k+1}, u_j \rangle - \sum_{i=1}^{k} \frac{\langle v_{k+1}, u_i \rangle}{\langle u_i, u_i \rangle} \langle u_i, u_j \rangle.
\end{align*}
By the inductive hypothesis, $\{u_1, \dots, u_k\}$ is orthogonal, so $\langle u_i, u_j \rangle = 0$ when $i \neq j$. Only the $i = j$ term survives:
\begin{align*}
\langle u_{k+1}, u_j \rangle
&= \langle v_{k+1}, u_j \rangle - \frac{\langle v_{k+1}, u_j \rangle}{\langle u_j, u_j \rangle} \langle u_j, u_j \rangle \\
&= \langle v_{k+1}, u_j \rangle - \langle v_{k+1}, u_j \rangle = 0.
\end{align*}
So $u_{k+1}$ is orthogonal to each of $u_1, \dots, u_k$.

\textbf{Non-zero:} By the inductive hypothesis, $\operatorname{span}\{u_1, \dots, u_k\} = \operatorname{span}\{v_1, \dots, v_k\}$. The formula for $u_{k+1}$ is $v_{k+1}$ minus a linear combination of $u_1, \dots, u_k$. If $u_{k+1} = 0$, this would mean $v_{k+1} \in \operatorname{span}\{u_1, \dots, u_k\} = \operatorname{span}\{v_1, \dots, v_k\}$. Since the $v_i$ are assumed to be linearly independent (they are non-zero vectors being fed into the process from a basis or independent set), this is a contradiction. Hence $u_{k+1} \neq 0$.

\textbf{Span:} Since $u_{k+1} = v_{k+1} - \sum_{i=1}^{k} r_i u_i$, we can write $v_{k+1} = u_{k+1} + \sum_{i=1}^{k} r_i u_i$, so $v_{k+1} \in \operatorname{span}\{u_1, \dots, u_{k+1}\}$. Conversely, $u_{k+1}$ is a linear combination of $v_{k+1}$ and $u_1, \dots, u_k$, and by the inductive hypothesis each $u_i \in \operatorname{span}\{v_1, \dots, v_k\}$, so $u_{k+1} \in \operatorname{span}\{v_1, \dots, v_{k+1}\}$.

Therefore
\begin{align*}
\operatorname{span}\{u_1, \dots, u_{k+1}\} = \operatorname{span}\{v_1, \dots, v_{k+1}\}.
\end{align*}

By induction, the result holds for all $k \ge 1$.

\hfill $\square$
}}

\bigskip

\clearpage
\hfill \break  $~\!\!$ \dotfill \medskip \\
\textbf{Question.}\\
Let $V$ be an inner product space, let $O = \{u_1, \dots, u_n\}$ be an orthonormal set, and let $S = \operatorname{span}(O)$. Define
\[
\hat{v} = \sum_{i=1}^{n} \langle v, u_i \rangle\, u_i.
\]
Show that if $\langle v - \hat{v}, u_i \rangle = 0$ for each $u_i \in O$, then $v - \hat{v} \in S^\perp$.

\noindent\\
\textbf{Solution.}\\
\quad\\
\noindent\fbox{ \parbox{\textwidth}{
We need to show that $\langle v - \hat{v}, s \rangle = 0$ for all $s \in S$.

Let $s \in S$. Since $O = \{u_1, \dots, u_n\}$ is a basis for $S$, there exist scalars $c_1, \dots, c_n \in F$ such that
\begin{align*}
s = c_1 u_1 + c_2 u_2 + \cdots + c_n u_n.
\end{align*}

Now compute:
\begin{align*}
\langle v - \hat{v}, s \rangle
&= \langle v - \hat{v},\; c_1 u_1 + \cdots + c_n u_n \rangle \\
&= c_1 \langle v - \hat{v}, u_1 \rangle + c_2 \langle v - \hat{v}, u_2 \rangle + \cdots + c_n \langle v - \hat{v}, u_n \rangle.
\end{align*}
Here we used linearity of the inner product in the second variable (conjugate-linearity gives $\overline{c_i}$ in the complex case; the argument is the same since each term is zero).

Wait---in the complex case, $\langle v - \hat{v}, c_i u_i \rangle = \overline{c_i}\, \langle v - \hat{v}, u_i \rangle$. But by hypothesis, $\langle v - \hat{v}, u_i \rangle = 0$ for each $i$. Therefore each term equals $\overline{c_i} \cdot 0 = 0$.

Hence
\begin{align*}
\langle v - \hat{v}, s \rangle = \overline{c_1} \cdot 0 + \overline{c_2} \cdot 0 + \cdots + \overline{c_n} \cdot 0 = 0.
\end{align*}

Since this holds for all $s \in S$, we conclude $v - \hat{v} \in S^\perp$.

\hfill $\square$
}}

\end{document}